
\documentclass{article}
\usepackage{colortbl}
\usepackage{makecell}
\usepackage{multirow}
\usepackage{supertabular}

\begin{document}

\newcounter{utterance}

\centering \large Interaction Transcript for game `cladder', experiment `full\_v1.5\_default', episode 5310 with selfrefine{-}Qwen3.5{-}27B.
\vspace{24pt}

{ \footnotesize  \setcounter{utterance}{1}
\setlength{\tabcolsep}{0pt}
\begin{supertabular}{c@{$\;$}|p{.15\linewidth}@{}p{.15\linewidth}p{.15\linewidth}p{.15\linewidth}p{.15\linewidth}p{.15\linewidth}}
    \# & \multicolumn{2}{c}{Player} && \multicolumn{2}{c}{Game Master} \\
    \hline

    \theutterance \stepcounter{utterance}  
    & & & \multicolumn{4}{p{0.6\linewidth}}{
        \cellcolor[rgb]{0.9,0.9,0.9}{
            \makecell[{{p{\linewidth}}}]{
                \texttt{\tiny{[P1$\langle$GM]}}
                \texttt{You are an expert in causal inference. The following question is not a typical commonsense query, but rather a meticulously designed question created by a professor specializing in causal inference, intended to assess the students' mastery of the course content.} \\
\texttt{Imagine a self{-}contained, hypothetical world with only the following conditions, and without any unmentioned factors or causal relationships: Liking spicy food has a direct effect on occupation and salary. Skill has a direct effect on occupation and salary. Occupation has a direct effect on salary. Skill is unobserved. Method 1: We look at how liking spicy food correlates with salary case by case according to occupation. Method 2: We look directly at how liking spicy food correlates with salary in general. To understand how liking spicy food affects salary, is it more correct to use the Method 1 than Method 2?} \\
\texttt{Start your answer with 'Yes' or 'No', followed by additional reasoning or evidence to support your explanation.} \\
            }
        }
    }
     \\ \\

    \theutterance \stepcounter{utterance}  
    & \multicolumn{4}{p{0.6\linewidth}}{
        \cellcolor[rgb]{0.561,0.737,0.561}{
            \makecell[{{p{\linewidth}}}]{
                \texttt{\tiny{[P1$\rangle$GM]}}
                \texttt{No, in this hypothetical world, **Skill** acts as a **confounder** because it has a direct causal effect on both the treatment (Liking spicy food) and the outcome (Salary), and it is unobserved. Here is the causal breakdown: 1. **The Confounding Path**: Skill \$\textbackslash{}rightarrow\$ Liking spicy food and Skill \$\textbackslash{}rightarrow\$ Salary. This creates a non{-}causal backdoor path between Liking spicy food and Salary. 2. **Method 2 (General Correlation)**: This estimates the total association, which includes the direct causal effect of Liking spicy food on Salary *plus* the spurious correlation induced by the unobserved Skill. Since Skill is unobserved, we cannot adjust for it, so Method 2 yields a biased estimate of the causal effect. 3. **Method 1 (Stratification by Occupation)**: This attempts to control for **Occupation**. However, in this specific causal graph, Occupation is a **mediator** (Liking spicy food \$\textbackslash{}rightarrow\$ Occupation \$\textbackslash{}rightarrow\$ Salary) and also a descendant of the confounder (Skill \$\textbackslash{}rightarrow\$ Occupation). If we condition on a mediator (Occupation), we block the indirect causal path through which Liking spicy food affects Salary. This results in estimating only the *direct* effect, not the *total* effect. More critically, because Occupation is a descendant of the unobserved confounder (Skill), conditioning on Occupation can induce **collider bias** (or selection bias) if there are other paths, or simply fail to remove the confounding bias from Skill. Since Skill affects both Liking spicy food and Occupation, and we are conditioning on Occupation, we do not block the backdoor path from Skill to Salary. The correlation between Liking spicy food and Salary within each occupation stratum will still be confounded by Skill. Therefore, neither method perfectly recovers the causal effect without observing Skill. However, Method 1 is **not more correct** than Method 2 for understanding the *total* effect of Liking spicy food on Salary. Method 1 introduces bias by blocking the causal pathway through Occupation (if the goal is total effect) and fails to eliminate the confounding bias from the unobserved Skill. In fact, conditioning on a variable (Occupation) that is influenced by an unobserved confounder (Skill) often exacerbates bias or leads to incorrect conclusions about the relationship. Without observing Skill, the causal effect is fundamentally unidentifiable using these methods alone, but Method 1 is structurally flawed for estimating the total effect due to the mediation issue.} \\
            }
        }
    }
    & & \\ \\

    \theutterance \stepcounter{utterance}  
    & & & \multicolumn{2}{p{0.3\linewidth}}{
        \cellcolor[rgb]{0.9,0.9,0.9}{
            \makecell[{{p{\linewidth}}}]{
                \texttt{\tiny{[GM$|$GM]}}
                \texttt{No} \\
            }
        }
    }
    & & \\ \\

    \theutterance \stepcounter{utterance}  
    & & & \multicolumn{2}{p{0.3\linewidth}}{
        \cellcolor[rgb]{0.9,0.9,0.9}{
            \makecell[{{p{\linewidth}}}]{
                \texttt{\tiny{[GM$|$GM]}}
                \texttt{no} \\
            }
        }
    }
    & & \\ \\

    \theutterance \stepcounter{utterance}  
    & & & \multicolumn{2}{p{0.3\linewidth}}{
        \cellcolor[rgb]{0.9,0.9,0.9}{
            \makecell[{{p{\linewidth}}}]{
                \texttt{\tiny{[GM$|$GM]}}
                \texttt{game\_result = WIN} \\
            }
        }
    }
    & & \\ \\

\end{supertabular}
}

\end{document}
